\documentclass[12pt]{article}
\usepackage{amsmath,amsthm,amssymb}
\usepackage{enumitem}
\title{Abstract Algebra - Homework 2}
\author{Simon Gustafsson}
\date{}

\begin{document}
\maketitle

\section*{Problem 1}

Write \(|G| = p^{\alpha}m\) with \(p\nmid m\) and \(|P|=p^{\alpha}\).  
Because \(N\trianglelefteq G\), Lagrange gives \(|N|=p^{\beta}d\) where \(0\le\beta\le\alpha\) and \(p\nmid d\). First consider \(N\cap P\).  Since \(N\cap P\le P\), we have
\[
|N\cap P|\mid |P|=p^{\alpha}\quad\Longrightarrow\quad|N\cap P|=p^{\gamma}
\]
for some \(0\le\gamma\le\beta\); thus \(N\cap P\) is a \(p\)-subgroup of \(N\).

Next form \(NP\le G\).  Using \(|NP|=\dfrac{|N||P|}{|N\cap P|}\) we compute
\[
|NP|=\frac{p^{\beta}d\;p^{\alpha}}{p^{\gamma}}=p^{\alpha+\beta-\gamma}d,
\qquad
[NP:P]=\frac{|NP|}{|P|}=p^{\beta-\gamma}d .
\]
Because \(P\le NP\le G\), the index \([NP:P]\) divides \([G:P]=m\); since \(p\nmid m\) we must have \(p^{\beta-\gamma}=1\).  Hence \(\beta=\gamma\) and \( |N\cap P|=p^{\beta}\).  This shows \(N\cap P\) attains the maximal power of \(p\) dividing \(|N|\), so \(N\cap P\) is a Sylow \(p\)-subgroup of \(N\).

For the quotient \(G/N\) we have
\[
|G/N|=\frac{|G|}{|N|}=\frac{p^{\alpha}m}{p^{\beta}d}=p^{\alpha-\beta}\frac{m}{d},
\]
so any Sylow \(p\)-subgroup of \(G/N\) has order \(p^{\alpha-\beta}\).
Since \(N\le NP\), the subgroup \(NP/N\le G/N\) satisfies
\[
|NP/N|=\frac{|NP|}{|N|}=\frac{p^{\alpha}d}{p^{\beta}d}=p^{\alpha-\beta},
\]
hence \(NP/N\) is a Sylow \(p\)-subgroup of \(G/N\).

Thus \(N\cap P\) is Sylow-\(p\) in \(N\), and \(NP/N\) is Sylow-\(p\) in \(G/N\).

\section*{Problem 2}
\begin{enumerate}[label=(\arabic*)]

\item
Let \(N=\ker f\).  For any subgroup \(S\le G\) one has
\[
f^{-1}\!\bigl(f(S)\bigr)=NS,
\]
because every element mapping into \(f(S)\) differs from an element of \(S\) by an element of \(N\).

Assume \(f(P)=f(Q)\).  Then
\[
P,\;Q\;\subseteq\;f^{-1}\!\bigl(f(P)\bigr)=NP ,
\]
so \(P\) and \(Q\) are both Sylow \(p\)-subgroups of the same subgroup \(NP\).
Sylow's conjugacy theorem applied inside \(NP\) gives an element
\(x\in NP\subseteq G\) with \(xPx^{-1}=Q\).
Hence \(P\) and \(Q\) are conjugate in \(G\).

\item
For the inclusion \(f\!\bigl(N_G(P)\bigr)\subseteq
N_H\!\bigl(f(P)\bigr)\) let \(g\in N_G(P)\); then \(gPg^{-1}=P\) and
\[
f(g)f(P)f(g)^{-1}=f(gPg^{-1})=f(P),
\]
so \(f(g)\in N_H(f(P))\).

Conversely, let \(h\in N_H(f(P))\).  Choose \(g\in G\) with \(f(g)=h\).
Since
\(f(gPg^{-1})=f(g)f(P)f(g)^{-1}=f(P)\),
part (1) implies \(gPg^{-1}\) is conjugate to \(P\) in \(G\):
there exists \(x\in G\) with \(xgPg^{-1}x^{-1}=P\).
Put \(y=xg\); then \(y\in N_G(P)\) and
\(h=f(g)=f(x)^{-1}f(y)\in f\!\bigl(N_G(P)\bigr)\).
Hence \(N_H(f(P))\subseteq f\!\bigl(N_G(P)\bigr)\).

Therefore \(f\!\bigl(N_G(P)\bigr)=N_H\!\bigl(f(P)\bigr)\).

\item
If \(p\nmid |H|\) then \(|f(P)|\mid|H|\) and \(|f(P)|\) is a power of \(p\); hence \(|f(P)|=1\) and \(f(P)=\{e\}\subseteq\ker f\).
From part (2) we have \(f\!\bigl(N_G(P)\bigr)=N_H(f(P))=H\),
because the normaliser of the trivial subgroup is all of \(H\).
\end{enumerate}

\section*{Problem 3}
\begin{enumerate}[label=(\arabic*)]

\item
(a)  By Sylow’s theorems, the number \(n_{5}\) of 5-Sylow subgroups satisfies  
\(n_{5}\equiv 1 \pmod 5\) and \(n_{5}\mid |G|/5=15/5=3\), so \(n_{5}=1\).  
A unique Sylow subgroup is normal, hence \(G\) has a normal subgroup of order \(5\).

\smallskip
(b)  Let \(P_{3}\) and \(P_{5}\) be the Sylow 3- and 5-subgroups of \(G\) (orders \(3\) and \(5\)).  
Because \(|P_{3}||P_{5}|=15=|G|\) and \(\gcd(3,5)=1\) we have \(P_{3}\cap P_{5}=\{e\}\).  
Part (a) gives \(P_{5}\trianglelefteq G\), and groups of prime order are cyclic, so
\[
G=P_{3}P_{5}\cong P_{3}\times P_{5}\cong\mathbb Z_{3}\times\mathbb Z_{5}\cong\mathbb Z_{15},
\]
hence \(G\) is cyclic.

\item
Again by Sylow’s theorems, \(n_{p}\equiv 1\pmod p\) and \(n_{p}\mid q\), so \(n_{p}=1\) or \(q\).  
Likewise \(n_{q}\equiv 1\pmod q\) and \(n_{q}\mid p\), so \(n_{q}=1\) or \(p\).  
If \(n_{q}=p\) then \(p\equiv 1\pmod q\), contradicting the hypothesis \(p\not\equiv 1\pmod q\); hence \(n_{q}=1\) and the \(q\)-Sylow subgroup \(Q\) is normal.

Both Sylow subgroups \(P,Q\) are cyclic of prime order, \(P\cap Q=\{e\}\) and \(|P||Q|=pq=|G|\), so
\[
G\cong P\times Q\cong\mathbb Z_{p}\times\mathbb Z_{q}\cong\mathbb Z_{pq},
\]
and \(G\) is cyclic.

\item
(a)  For \(|G|=12=2^{2}\!\cdot 3\).  
Sylow gives \(n_{3}\equiv 1\pmod 3,\; n_{3}\mid 4\Rightarrow n_{3}=1\text{ or }4\);  
\(n_{2}\equiv 1\pmod 2,\; n_{2}\mid 3\Rightarrow n_{2}=1\text{ or }3\).

If \(n_{3}=1\) \emph{or} \(n_{2}=1\) we already have a normal Sylow subgroup.  
Assume instead \(n_{3}=4\) and \(n_{2}=3\).  
Each 3-Sylow subgroup contributes \(2\) non-identity elements, so \(4\) of them give \(8\) such elements.  
Each 2-Sylow subgroup (of order 4) contributes \(3\) non-identity elements; \(3\) of them would add \(9\) more, totalling \(17>12\), impossible.  
Therefore at least one Sylow subgroup is unique and normal, so \(G\) always has a normal subgroup.

\smallskip
(b)  A non-cyclic group of order 12 is the alternating group \(A_{4}\).  
A cyclic group of order 12 would contain an element of order 12, but in \(A_{4}\) every element has order \(1,2,\) or \(3\); hence \(A_{4}\) is not cyclic.

\end{enumerate}
\section*{Problem 4}
First, note that both \( a \) and \( b \) are products of two disjoint 4-cycles. Therefore,
\[
|a| = |b| = \mathrm{lcm}(4, 4) = 4.
\]

Since their cycles overlap, they do not act on disjoint sets, so we compute explicitly:
\[
ab = (1234)(5678)(1537)(2648), \quad ba = (1537)(2648)(1234)(5678).
\]
We compute
\[
ab=(17)(28)(36)(45)=ba,
\]
so \(a\) and \(b\) commute and \(H\) is abelian.

To determine the order of \( H \), observe that \( a \) has order 4, and \( b \) also has order 4. We compute \( a^2 \), \( b^2 \), and their combinations to find the number of distinct elements generated. It turns out that:
\[
H = \langle a, b \rangle = \{e, a, a^2, a^3, b, ab, a^2b, a^3b\}
\]
which contains 8 distinct elements. So \( |H| = 8 \).

Since \( H \) is a finite abelian group of order 8, by the classification theorem for finite abelian groups, it must be isomorphic to one of:
\[
\mathbb{Z}_8, \quad \mathbb{Z}_4 \times \mathbb{Z}_2, \quad \mathbb{Z}_2 \times \mathbb{Z}_2 \times \mathbb{Z}_2.
\]

We note that \( H \) contains elements of order 4 (e.g., \( a \)), but no element of order 8, so it cannot be cyclic. It also contains an element of order 4 and another of order 2 (e.g., \( b^2 \)), so it is not the elementary abelian group \( \mathbb{Z}_2 \times \mathbb{Z}_2 \times \mathbb{Z}_2 \), which has all elements of order 2.

Therefore, the group must be:
\[
H \cong \mathbb{Z}_4 \times \mathbb{Z}_2.
\]
\section*{Problem 5}
\begin{enumerate}[label=(\arabic*)] 

\item Assume $G/Z$ is cyclic, so there exists $g \in G$ such that every element of $G/Z$ is of the form $g^nZ$ for some $n \in \mathbb{Z}$. Then for any $x \in G$, there exists $n \in \mathbb{Z}$ and $z \in Z$ such that:
\[
x = g^n z.
\]

Now let $x = g^n z$ and $y = g^m z'$, with $z, z' \in Z$. Then:
\[
xy = g^n z \cdot g^m z' = g^{n+m} zz', \quad \text{and} \quad yx = g^m z' \cdot g^n z = g^{m+n} z'z.
\]

Since $Z \subseteq Z(G)$, we have $zz' = z'z$, so:
\[
xy = yx.
\]

Thus, any two elements of $G$ commute, so $G$ is abelian.

\item Since $G$ is a $p$-group, its center $Z(G)$ is nontrivial. Hence, $|Z(G)| = p$ or $p^2$. 
If $|Z(G)| = p^2$, then $Z(G) = G$, so $G$ is abelian. 
If $|Z(G)| = p$, then the quotient group $G/Z(G)$ has order $p$ and is cyclic. By previous section, if $G/Z(G)$ is cyclic and $Z(G)$ is central, then $G$ is abelian. 
Therefore, in both cases, $G$ is abelian.

\item Let $p$ be a prime. Up to isomorphism, there are exactly two groups of order $p^2$:
\[
\mathbb{Z}_{p^2}, \quad \mathbb{Z}_p \times \mathbb{Z}_p.
\]

These are the only abelian groups of order $p^2$, and by previous sectio, all groups of this order are abelian. Therefore, this list is complete up to isomorphism.
\end{enumerate}
\section*{Problem 6}
\begin{enumerate}[label=(\arabic*)] 
\item Since $A \cong \mathbb{Z}_n$, any homomorphism $\chi \in \widehat{A}$ is determined by the image of a generator, and this image must be an $n$-th root of unity. Therefore, $\widehat{A} \cong \mu_n$, the group of $n$-th roots of unity, which is cyclic of order $n$. 
Hence, $\widehat{A}$ is cyclic of order $n$.

\item By the structure theorem for finite abelian groups, $A$ is isomorphic to a product of cyclic groups:
\[
A \cong \mathbb{Z}_{n_1} \times \cdots \times \mathbb{Z}_{n_k}.
\]
Then,
\[
\widehat{A} \cong \widehat{\mathbb{Z}_{n_1}} \times \cdots \times \widehat{\mathbb{Z}_{n_k}},
\]
and each $\widehat{\mathbb{Z}_{n_i}} \cong \mathbb{Z}_{n_i}$. Therefore, $\widehat{A} \cong A$, so $|\widehat{A}| = |A| = n$.
\end{enumerate}


\end{document}